\section*{Chapitre \ref{ch:b3:structural-analysis} --- Analyse structurale}

\begin{solution}{exo:b3:structural-analysis:1}
Le \omterm{def:b3:structural-analysis:stereo}{pôle} a un azimut $40 + 180 = \ang{220}$ et un plongement $90 - 30 = \ang{60}$. Le
plan a un pendage de $90 - 20 = \ang{70}$ vers $315 - 180 = \ang{135}$~: $135/70$.
\end{solution}

\begin{solution}{exo:b3:structural-analysis:2}
$\theta = \ang{60}$. \omterm{def:b3:structural-analysis:stereo}{Canevas équiangle}~: $10\tan\ang{30} = \qty{5.77}{cm}$~;
\omterm{def:b3:structural-analysis:stereo}{projection équivalente}~: $10\sqrt2\sin\ang{30} = \qty{7.07}{cm}$.
\end{solution}

\begin{solution}{exo:b3:structural-analysis:3}
$\vec n_1 = (\sin\ang{30}, 0, \cos\ang{30})$, $\vec n_2 = (-\sin\ang{30}, 0,
\cos\ang{30})$, $\vec n_1\cdot\vec n_2 = -0.25 + 0.75 = 0.5$~: les \omterm{def:b3:structural-analysis:stereo}{pôles} sont
distants de \ang{60}, l'angle entre flancs vaut \ang{120}, un pli ouvert.
\end{solution}

\begin{solution}{exo:b3:structural-analysis:4}
$\vec d = (\cos\ang{60}, 0, -\sin\ang{60}) = (0.5, 0, -0.866)$ et $\vec u =
0.707\,\vec s + 0.707\,\vec d = (0.354, 0.707, -0.612)$~: azimut
$\operatorname{atan2}(0.354, 0.707) = \ang{26.6}$, plongement $\arcsin0.612 =
\ang{37.8}$. Le pitch se situe entre \ang{-90} (normal) et 0 (senestre)~: une
faille normale oblique à composante senestre.
\end{solution}

\begin{solution}{exo:b3:structural-analysis:5}
\begin{gather*}
  \vec n_1 = (0.5567, 0.3214, 0.7660),\qquad \vec n_2 = (-0.4698, -0.1710, 0.8660),\\
  \vec n_1\times\vec n_2 = (0.4093, -0.8420, 0.0558),
\end{gather*}
{\sloppy de longueur $0.9379$. Son opposé, $-\vec n_1\times\vec n_2$, pointe vers le bas~:
azimut \ang{334.1}, plongement $\arcsin(0.0558/0.9379) = \ang{3.4}$. $\vec n_1\cdot\vec
n_2 = 0.3469$, \ang{69.7} entre les normales, angle entre flancs \ang{110.3}~:
un pli ouvert.\par}
\end{solution}

\begin{solution}{exo:b3:structural-analysis:6}
Un plongement supérieur à \ang{60} signifie $\theta < \ang{30}$. \omterm{def:b3:structural-analysis:stereo}{Projection équivalente}~: $r^2 =
2\sin^2\ang{15} = 0.134$, \qty{13.4}{\percent} du canevas~; avec le \omterm{def:b3:structural-analysis:stereo}{canevas équiangle}~:
$\tan^2\ang{15} = 0.072$, \qty{7.2}{\percent}. Sur la sphère, la calotte représente
$1 - \cos\ang{30} = 0.134$ de l'hémisphère~: seule la \omterm{def:b3:structural-analysis:stereo}{projection équivalente}
la restitue fidèlement, donc c'est sur elle que l'on compte les \omterm{def:b3:structural-analysis:stereo}{pôles}.
\end{solution}

\begin{solution}{exo:b3:structural-analysis:7}
\begin{gather*}
  T_{xx} = \frac{2\sin^2\ang{20}}3 = 0.078,\qquad
  T_{zz} = \frac{2\cos^2\ang{20} + 1}3 = 0.922,\\
  T_{xz} = -\frac{\sin\ang{20}\cos\ang{20} - \sin\ang{20}\cos\ang{20}}3 = 0,
\end{gather*}
et tous les termes en $y$ sont nuls. Les valeurs propres sont $0.922$ (verticale),
$0.078$ (est) et $0$ (nord)~: l'\omterm{def:b3:structural-analysis:pi}{axe bêta} est horizontal, nord--sud.
Les trois \omterm{def:b3:structural-analysis:stereo}{pôles} sont dans un même plan, donc aucun n'a de composante selon sa normale,
et $\vec u^{\mathsf T}T\vec u = 0$ pour $\vec u$ dirigé selon ce plan.
\end{solution}

\begin{solution}{exo:b3:structural-analysis:8}
$a = 2$, $b = 1.5$, $k = 1/0.5 = 2$~: constriction. $Z = 1/(1.8\times1.1) =
0.505$, $a = 1.636$, $b = 2.178$, $k = 0.636/1.178 = 0.54$~: aplatissement.
\end{solution}

\begin{solution}{exo:b3:structural-analysis:9}
$S_1/S_2 = 2.25$ et $S_1S_2 = 1$~: $S_1 = 1.5$, $S_2 = 0.667$~; le grand axe
a été allongé de \qty{50}{\percent}. Le vide a pour demi-axes $1.5\times0.4 =
\qty{0.60}{mm}$ et $0.667\times0.4 = \qty{0.27}{mm}$.
\end{solution}

\begin{solution}{exo:b3:structural-analysis:10}
{\sloppy
\omterm{def:b3:structural-analysis:stereo}{Canevas équiangle}~: $1/(2\cos^2\ang{30}) = 0.667$ selon la direction radiale et
$\tan\ang{30}/\sin\ang{60} = 0.577/0.866 = 0.667$ en direction orthoradiale~: ils sont égaux et la projection est conforme.
\omterm{def:b3:structural-analysis:stereo}{Projection équivalente}~: $\mathrm dr/(R\,\mathrm d\theta) = (\sqrt2/2)\cos\ang{30} = 0.612$
et $\sqrt2\sin\ang{30}/\sin\ang{60} = 0.707/0.866 = 0.816$~; le produit vaut
$0.500$, soit le facteur d'aire du théorème. Un petit cercle devient un ovale d'axes
dans le rapport $0.816/0.612 = 1.33$, dont le grand axe est orienté autour du centre.\par}
\end{solution}

\begin{solution}{exo:b3:structural-analysis:11}
La faille est verticale, donc sa normale est horizontale et perpendiculaire à la
verticale $\sigma_2$~; d'après la partie (b) du théorème, le glissement est perpendiculaire à
$\sigma_2$, horizontal, pour tout $\Phi$. Avec $\sigma_1 = 1$ selon le nord,
$\sigma_3 = 0$ selon l'est~: $S\vec n = (0, -0.342, 0)$, $\vec n\cdot S\vec n =
0.117$, $\vec\tau = (0, -0.342, 0) - 0.117\,(0.940, -0.342, 0) = (-0.110,
-0.302, 0)$ et $-\vec\tau = (0.110, 0.302, 0)$, vers \ang{020}. Le bloc
du côté de $\vec n$, le côté est, se déplace vers le nord-nord-est le long de la
direction~: vu depuis le côté ouest, il se déplace vers la gauche, une faille senestre (pitch
0). $\sigma_2$ n'est pas intervenu dans le calcul.
\end{solution}

\begin{solution}{exo:b3:structural-analysis:12}
$S = (7.5 - 6.0)/1.0 = \qty{1.5}{km}$~; $h = 6.0/1.5 = \qty{4.0}{km}$ sous
l'horizon inférieur. La coupe mesurait $20 + 1.5 = \qty{21.5}{km}$ de long, et a varié
de $-1.5/21.5 = \qty{-7.0}{\percent}$.
\end{solution}

\begin{solution}{pb:b3:structural-analysis:1}
\textbf{1.} Le \omterm{def:b3:structural-analysis:stereo}{pôle} a pour azimut \ang{255} et pour plongement \ang{77}~; $\theta =
\ang{13}$ et $r = 10\sqrt2\sin\ang{6.5} = \qty{1.60}{cm}$ du centre,
vers \ang{255}.

\textbf{2.} $\vec n = (\sin\ang{13}\sin\ang{75}, \sin\ang{13}\cos\ang{75},
\cos\ang{13}) = (0.2173, 0.0582, 0.9744)$. Six vecteurs unitaires presque parallèles
s'additionnent presque à 6~: $|\vec S_E| = 5.9974$.

\textbf{3.} Flanc est~: $\bar\alpha = \operatorname{atan2}(1.2583, 0.2017) =
\ang{80.9}$ et $\bar\delta = \arccos(5.8604/5.9974) = \ang{12.3}$~; $\bar R
= 5.9974/6 = 0.99956$ et la dispersion vaut $\sqrt{2\times0.00044} =
\qty{0.0297}{rad}$, soit \ang{1.7}. Flanc ouest~: $\bar\alpha =
\operatorname{atan2}(-1.2445, 0.1886) = \ang{278.6}$ et $\bar\delta =
\arccos(5.8636/5.9972) = \ang{12.1}$~; $\bar R = 0.99953$, dispersion \ang{1.8}.

\textbf{4.}
\[
  \vec S_E\times\vec S_W = (0.0774, -14.6714, 0.4883),
\]
de longueur $14.680$. Son opposé pointe vers le bas~: l'azimut vaut $\operatorname{atan2}(-0.0774, 14.6714) = \ang{359.7}$ et le
plongement $\arcsin(0.4883/14.680) = \ang{1.9}$.

\textbf{5.} Le cosinus de l'angle entre les normales moyennes vaut
$32.835/35.968 = 0.9129$, le produit scalaire des sommes divisé par leurs
longueurs~: \ang{24.1} entre les normales, soit un angle entre flancs de \ang{155.9}~:
un pli ouvert, voisin d'un pli doux.

\textbf{6.}
\begin{gather*}
  \vec m_E = (0.2098, 0.0336, 0.9772),\qquad \vec m_W = (-0.2075, 0.0314, 0.9777),\\
  \vec m_E - \vec m_W = (0.4173, 0.0022, -0.0005):
\end{gather*}
une normale horizontale vers \ang{089.7}. Le plan axial est vertical, de direction
\ang{359.7}, pratiquement nord--sud~: le pli est droit.

\textbf{7.} La ceinture est le plan dont le \omterm{def:b3:structural-analysis:stereo}{pôle} est l'\omterm{def:b3:structural-analysis:pi}{axe bêta}~: sens du pendage
\ang{179.7}, pendage $90 - 1.9 = \ang{88.1}$, presque vertical, dans le
plan est--ouest. Chaque flanc contient l'axe de pli, donc chaque normale est perpendiculaire
à lui et se situe dans le plan qui lui est perpendiculaire.

\textbf{8.} $0.04400 + 0.00144 + 0.95456 = 1.00000$. $\vec b^{\mathsf T}T\vec
b = 0.00038$, la plus petite valeur propre~: $\vec b$ est, à l'arrondi près, le
vecteur propre $\vec e_3$, et la ceinture passant par les deux \omterm{def:b3:structural-analysis:stereo}{pôles} moyens est la
meilleure ceinture d'ajustement parmi les douze.

\textbf{9.} $\arcsin\sqrt{0.00038} = \arcsin0.0195 = \ang{1.1}$.

\textbf{10.} Le long de la ceinture, les \omterm{def:b3:structural-analysis:stereo}{pôles} se situent à environ \ang{12.2} de la verticale,
de part et d'autre~; $\sin^2\ang{12.2} = 0.045$, le carré moyen de leurs
composantes est--ouest, soit $\lambda_2$.

\textbf{11.} Oui~: les \omterm{def:b3:structural-analysis:stereo}{pôles} s'écartent de la ceinture de \ang{1.1} en moyenne, pas plus
que la dispersion de \ang{1.7} des mesures d'un seul flanc autour de leur moyenne.

\textbf{12.} $\ang{1.1}/\sqrt{12} = \ang{0.3}$~: le plongement vaut \ang{1.9} $\pm$
\ang{0.3}.

\textbf{13.} Un pitch de \ang{-90}~: une faille normale pure, à glissement de pendage. Anderson~:
$\sigma_1$ verticale, $\sigma_3$ horizontale selon le sens du pendage,
\ang{120}--\ang{300}, $\sigma_2$ horizontale selon la direction,
\ang{030}--\ang{210}.

\textbf{14.} La normale de $120/60$ se situe dans le plan vertical d'azimut
\ang{120}, qui contient $\vec e_1$ et $\vec e_3$~: $n_2 = 0$. D'après la question~(b), le
glissement est perpendiculaire à $\sigma_2$, qui suit la direction, quel que soit $\Phi$~: il est
à glissement de pendage. La faille ne dit rien de $\Phi$~; dans la formule, $n_2$ et
$\vec e_2\cdot\vec w$, donc numérateur et dénominateur s'annulent.

\textbf{15.} $\vec n = (\sin\ang{55}\sin\ang{150}, \sin\ang{55}\cos\ang{150},
\cos\ang{55}) = (0.4096, -0.7094, 0.5736)$; $n_1 = 0.574$; $n_2 = 0.5\times
0.4096 + 0.866\times(-0.7094) = -0.410$.

\textbf{16.} $\vec e_1\cdot\vec w = 0.198$ et $\vec e_2\cdot\vec w = 0.5\times
0.771 + 0.866\times0.605 = 0.909$, donc $\Phi = -(0.574\times0.198)/(-0.410\times
0.909) = 0.1137/0.3727 = 0.305$.

\textbf{17.} $(0.305 + 0.294 + 0.301 + 0.306)/4 = 0.30$. Les pitchs ont été
arrondis au degré près~; un degré de pitch déplace $\Phi$ d'environ 0.02 sur ces
failles, si bien que l'arrondi au demi-degré explique la dispersion.

\textbf{18.} $\sigma_2 - \sigma_3 = 0.3(\sigma_1 - \sigma_3)$~: $\sigma_2$ est
plus proche de $\sigma_3$ que de $\sigma_1$. Les deux contraintes horizontales sont nettement
inférieures à la verticale~; la région a été étirée surtout vers \ang{120}, un
peu selon \ang{030}.

\textbf{19.} Avec $\vec n = (0.750, -0.433, 0.500)$ et le glissement $\vec u =
(0.433, -0.250, -0.866)$, $(\vec n - \vec u)/\sqrt2$ plonge de \ang{75} vers
\ang{300}~: P~; $(\vec n + \vec u)/\sqrt2$ plonge de \ang{15} vers \ang{120}~:
T. Ils sont dans le même plan vertical que $\sigma_1$ et $\sigma_3$, mais
décalés de \ang{15} par rapport à eux, car la faille fait \ang{30}, et non \ang{45}, avec
$\sigma_1$.

\textbf{20.} $\tan a = \tan\ang{12.3}\times\cos(\ang{90} - \ang{80.9}) =
0.2180\times0.9874 = 0.2153$: $a = \ang{12.1}$.

\textbf{21.} $L_0 = 8000/\cos\ang{12} = \qty{8178.7}{m}$, $S =
\qty{178.7}{m}$, une variation relative de $\cos\ang{12} - 1 =
\qty{-2.2}{\percent}$.

\textbf{22.} $4000\tan\ang{12} = \qty{850.2}{m}$; $A = \tfrac12\times8000\times
850.2 = \qty{3.40e6}{m^2}$.

\textbf{23.} $h = \num{3.40e6}/178.7 = \qty{19.0}{km}$. Un décollement à la
base de la couverture donnerait $h$ de quelques centaines de mètres et une aire
excédentaire environ cinquante fois plus petite~: les plis ne sont pas décollés à cet endroit. Soit le
socle s'est plissé avec la couverture, soit les synclinaux se sont enfoncés sous le
niveau initial, ce qui réduit $A$~; la coupe seule ne permet pas de trancher.

\textbf{24.} Rejet horizontal $200/\tan\ang{60} = \qty{115.5}{m}$, glissement $200/\sin\ang{60}
= \qty{230.9}{m}$.

\textbf{25.} $\tan a = \tan\ang{60}\cos\ang{30} = 1.500$~: $a = \ang{56.3}$, et
le rejet horizontal apparent vaut $200/1.5 = \qty{133.3}{m}$. Le glissement de la faille est
selon le pendage, vers \ang{120}~: la coupe selon \ang{120} montre les vraies
longueurs et c'est celle qu'il faut restaurer.

\textbf{26.} La faille recoupe les plis, elle est donc plus jeune et on la restaure en premier,
en faisant remonter le toit le long de la faille de \qty{230.9}{m}~; les plis sont
alors dépliés. Les \omterm{def:b3:structural-analysis:restoration}{lignes de référence} se placent sur les surfaces axiales des synclinaux, dans le
mur, là où les plis symétriques ne présentaient aucun glissement entre les couches.

\textbf{27.} L'axe du pli de Varden plonge
\textbf{\ang[text-series-to-math=true]{1.9}} vers le nord (azimut
\ang{359.7})~: le \ang{2} imposé à la région, retrouvé à la
précision de \ang{0.3} qu'autorisent les douze mesures.
\end{solution}
